\section{Justificación técnica y beneficios}
La implementación de un sistema de cableado estructurado en la Escuela Profesional
de Enfermería responde a la necesidad de contar con una infraestructura de red
estable, organizada y preparada para soportar el crecimiento tecnológico de la
institución.

Actualmente, las redes improvisadas o sin estandarización generan problemas de
rendimiento, dificultades de mantenimiento y limitaciones para futuras ampliaciones.
Por ello, se propone una solución basada en estándares internacionales que garantice
eficiencia, estabilidad y facilidad de administración.

Entre los principales beneficios del proyecto se encuentran:

\begin{itemize}
    \item Mayor velocidad y estabilidad en la transmisión de datos.
    
    \item Mejor organización y administración del cableado.
    
    \item Red preparada para futuras ampliaciones tecnológicas.
    
    \item Soporte para servicios de voz, datos y conectividad inalámbrica.
    
    \item Reducción de fallas y tiempos de mantenimiento.
    
    \item Mejora de los servicios académicos y administrativos.
\end{itemize}